\chapter{Cronograma Provisório}

O cronograma será ajustado após a confirmação da turma, da lotação docente, do componente curricular, do calendário escolar e das exigências institucionais e éticas. A Tabela~\ref{tab:cronograma} apresenta a ordem de trabalho prevista; ela não substitui o calendário acadêmico oficial nem as condições que ainda dependem da definição do campo.

\begin{table}[htbp]
\centering
\caption{Etapas e atividades do cronograma provisório da pesquisa}
\label{tab:cronograma}
\begin{tblr}{
  width = \textwidth,
  colspec = {X[1.2,l,m] X[3.8,l,m] X[3.0,l,m]},
  hline{1,Z} = {1.5pt},
  hline{2} = {1pt},
  rows = {abovesep=4pt,belowsep=4pt},
}
\textbf{Etapa} & \textbf{Atividades principais} & \textbf{Condição para conclusão} \\
1. Consolidação do projeto & Revisar problema, objetivos, referencial, matriz curricular, sequência e instrumentos. & Versão de qualificação coerente e submetida à apreciação do orientador. \\
2. Aperfeiçoamento do produto & Padronizar cenários sintéticos, adequar textos e controles do site à sequência, criar síntese comparativa e realizar verificações de acessibilidade e projeção. & Protótipo compatível com a aplicação planejada. \\
3. Preparação do campo & Confirmar turma, escola, componente curricular, infraestrutura, calendário e autorizações aplicáveis. & Campo definido e condições pedagógicas e éticas confirmadas. \\
4. Revisão dos instrumentos & Ajustar diagnóstico, desafio final, rubrica e registro docente ao perfil da turma. & Instrumentos revisados e prontos para uso. \\
5. Aplicação da intervenção & Realizar os oito encontros, produzir registros docentes e aplicar diagnóstico e desafio final. & Produções organizadas por códigos, sem identificação nominal no corpus de análise. \\
6. Organização e análise & Anonimizar, organizar e analisar respostas e registros; elaborar sínteses, frequências descritivas e excertos. & Categorias analíticas e resultados interpretados. \\
7. Redação da dissertação e do produto & Redigir resultados, discussão, considerações finais e versão final do produto educacional. & Versão submetida às revisões do orientador. \\
8. Revisão e entrega & Revisar normalização, ética, acessibilidade, anexos e exigências de depósito. & Trabalho e produto preparados para a etapa institucional correspondente. \\
\end{tblr}
\source{Elaborado pelo autor (2026).}
\end{table}
